\section{Environmental Disruptions, Evacuation, and Mobility Adaptation}
\label{sec:results_environment_disruptions}

Environmental disruptions are a useful test of whether an urban population model can distinguish a change in the environment from the behavioral response to that change. A location may become unavailable, but this state alone does not determine what the population should do. A routine visit may be cancelled or redirected, a work or school trip may remain unmet because its destination is not substitutable, and a resident exposed at home may seek protection in a capacity-constrained shelter. These responses require different policies even when they originate from the same type of environmental observation.

The experiments in this section extend the routine-mobility baselines from Section~\ref{sec:results_routine_mobility} along two complementary paths. The first path applies a time-varying water-level series to homes and activity, work, and school destinations, then evaluates trip suppression, destination substitution, and flood-aware commuting. The second path maps a static 5.30~m exposure estimate to the full-city population and evaluates evacuation toward shelters under capacity, response-rate, occupancy, distance, reallocation, and facility-failure variations. Together, the studies examine both disruption of ordinary movement and movement initiated specifically in response to exposure.

\subsection{Environmental-Disruption Plugins}
\label{sec:environment_disruption_plugins}

\subsubsection{Water-Level and Node-Availability Plugin}

The \textit{Water-Level Data Plugin} reads water levels indexed by simulation day and hour. Each target node contains a water-level threshold. At every step, the plugin compares the current level with that threshold and adds or removes a \texttt{flood} disable cause. Experiment configuration determines which node types are exposed: homes; marketplaces, restaurants, and pharmacies; work and school destinations; or a combination of these types. Because the environmental plugin changes availability rather than directly moving population, each mobility plugin remains responsible for its own response.

Popular Times V2 excludes disabled homes from traveler sourcing and suppresses a visit when its destination is unavailable. It can also anticipate whether a destination will become unavailable before the one-hour visit ends. The original L\'evy Walk configurations enable opt-in origin and destination availability checks, preventing routine movement to or from disabled nodes. L\'evy Walk V2 additionally records demand lost at disabled origins, unavailable or imminently unavailable destinations, and insufficiently available population. Its traceable commute lifecycle retains the original home so that temporary returns, blocked returns, and later repatriation can be audited.

\subsubsection{Destination-Adaptation Policy}

Popular Times V2 provides an optional nearest-enabled, same-type substitution policy. When an activity destination is unavailable, the original request is redirected to the geographically nearest enabled marketplace, restaurant, or pharmacy of the same type. The demand remains associated with the original destination and paired home, allowing the logger to distinguish the requested and receiving points of interest and to measure added point-of-interest displacement and receiving load.

This policy is deliberately limited to activity visits. Work and school attendance in the reported L\'evy Walk V2 scenarios uses suppression when the requested destination is unavailable. The distinction tests a modeling assumption: an occasional visit to a restaurant or pharmacy may be treated as substitutable, while a workplace or school represents a less interchangeable destination. Substitution is therefore a plugin-level policy rather than an automatic consequence of node availability.

\subsubsection{Shelter Plugin}

The \textit{Shelter Plugin} represents exposed residents and shelter use through a traceable \texttt{flooding\_status} characteristic with three values: \texttt{safe}, \texttt{in\_danger}, and \texttt{sheltered}. During initialization, the plugin maps census-sector exposure requests to the population available in each region. It then creates shelter nodes from the configured shelter file, loads their capacities and initial occupancies, and moves the initial sheltered population from exposed origins to those nodes.

Three actions define the evacuation lifecycle. The \texttt{move\_to\_shelters} action selects a configured share of exposed residents and invokes a direct L\'evy movement toward enabled shelter nodes. The \texttt{shelter\_population} action admits waiting residents without exceeding the receiving shelter's capacity; denied residents remain on its waitlist with an \texttt{in\_danger} status. Finally, \texttt{reallocate\_shelter\_waitlist} moves waiting residents from a full shelter to enabled shelters with free places. Reallocation can redistribute unused capacity, but it neither creates capacity nor changes the number of exposed residents.

The shelter logger records exposure, requests, arrivals, admissions, denials, reallocations, occupancy, utilization, unresolved demand, waiting-person-hours, origin--shelter flows, travel distance, and age- and occupation-group outcomes. It also retains an input audit covering region-name aliases, population shortfalls, capacity adjustments, loaded shelters, and configured shelter failures.

\subsection{Simulation Setup}
\label{sec:environment_disruption_setup}

\subsubsection{Mobility-Disruption Scenarios}

The mobility experiments use the reduced Porto Alegre representation with 13 regions, 435 census-sector home nodes, and a population of 131,925. Each census sector also has work, school, marketplace, restaurant, and pharmacy nodes. All scenarios cover 56 days with 24 hourly steps per day. Table~\ref{tbl:mobility_disruption_batches} summarizes the three completed batches used in this section.

\begin{table}[ht]
\centering
\footnotesize
\caption{Mobility-disruption and adaptation batches. The Stage 4 count includes the no-disruption references required for paired comparisons but excludes the separate 94-region baseline runs.}
\label{tbl:mobility_disruption_batches}
\begin{tabularx}{\linewidth}{@{}llrrX@{}}
\toprule
\textbf{Batch} & \textbf{Movement model} & \textbf{Scenarios} & \textbf{Runs} & \textbf{Configured conditions} \\
\midrule
Popular Times core
    & Original L\'evy off/on & 8 & 124
    & No exposure, activity destinations, homes, or both; deterministic L\'evy-off runs and seeds 0--29 with L\'evy enabled. \\
Activity adaptation
    & Original L\'evy off/on & 2 & 31
    & Homes and activity destinations exposed, with nearest-enabled same-type destination substitution. \\
Flood-aware commuting
    & L\'evy Walk V2 & 5 & 150
    & No exposure, destinations, homes, all node types, or all node types with activity-destination substitution; seeds 0--29. \\
\bottomrule
\end{tabularx}
\end{table}

The core Popular Times matrix exposes only homes and activity points of interest. Its destination condition therefore affects marketplaces, restaurants, and pharmacies. In the L\'evy Walk V2 matrix, the destination condition additionally affects work and school nodes, making it possible to measure activity fulfillment and commute fulfillment separately. All 305 runs used by these comparisons passed their population-conservation, demand-accounting, enabled-node movement, and plugin-specific lifecycle checks.

\subsubsection{Shelter Scenarios}

The shelter experiments use the full 94-region environment. It contains 598 base points of interest and 104 added shelter nodes, for 702 nodes in total, and represents 1,409,464 people. The 5.30~m exposure input identifies 186,522 people in 438 census sectors. Region-level population caps allow 182,936 of them to be mapped to available home populations, leaving an audited shortfall of 3,586. The shelter input provides 12,519 places and a configured initial occupancy of 10,963.

The spatial input audit contains 2,649 census-sector identifiers, of which 2,557 match the bundled geometry and 92 do not. These unmatched geometries limit the completeness of sector-level maps; exposure quantities are still resolved through their configured region mappings. Region names use normalized exact matching followed by a versioned alias table, and unresolved names stop execution rather than being matched fuzzily.

Each research scenario covers 14 days with 24 hourly steps. Evacuation requests occur at 07:00, initial admission at 08:00, optional global waitlist reallocation at 09:00, and admission of reallocated residents at 10:00. The reference response requests 1\% of the initially exposed population per day and uses a 1,000~m L\'evy distribution scale. Table~\ref{tbl:shelter_scenario_families} presents the screening catalog. Each of its 28 configurations was run with seeds 0--4, producing 140 completed or validly resumed runs. The conservation difference was zero at every recorded step in all 140 runs.

\begin{table}[ht]
\centering
\footnotesize
\caption{Fourteen-day shelter research catalog. Each configuration uses five seeds.}
\label{tbl:shelter_scenario_families}
\begin{tabularx}{\linewidth}{@{}lrlX@{}}
\toprule
\textbf{Family} & \textbf{Configurations} & \textbf{Values} & \textbf{Question} \\
\midrule
Reference and reallocation & 2 & Off; on & Does moving waitlists between shelters use spare capacity more efficiently? \\
Evacuation demand basis & 2 & Fixed initial; remaining demand & Is each day's request based on initial exposure or currently unresolved exposure? \\
Daily response rate & 4 & 0.5\%, 1\%, 2\%, 5\% & How quickly is evacuation demand introduced? \\
Exposure magnitude & 3 & 0.25, 0.5, 1.0 times & How does demand relative to fixed capacity affect protection? \\
Shelter capacity & 5 & 0.5, 1, 2, 5, 15 times & When does capacity cease to be the binding constraint? \\
Initial occupancy & 3 & Empty; half; configured & How does pre-existing occupancy affect waiting and travel? \\
L\'evy scale & 3 & 500, 1,000, 3,000~m & How does destination sampling affect travel distance? \\
Shelter failure & 6 & Three failures, reallocation off/on & How do failure topology and reallocation change outcomes? \\
\bottomrule
\end{tabularx}
\end{table}

Shelter coverage is the final \texttt{sheltered} population divided by mapped exposure. It includes the configured initial occupancy and therefore should not be read as the number of new admissions. Unresolved demand is the final \texttt{in\_danger} population. The peak waitlist and accumulated waiting-person-hours describe congestion after arrival, while mean and 95th-percentile distance describe the origin-to-shelter movements that occurred. Reported shelter values are five-seed means; the preserved analysis uses bootstrap intervals, although outcomes determined exactly by fixed exposure and capacity have no across-seed variation.

\subsection{Results}
\label{sec:environment_disruption_results}

\subsubsection{Popular Times Suppression and Adaptation}

Table~\ref{tbl:popular_times_disruption} reports the Popular Times results with the original L\'evy routine enabled. Destination exposure suppresses 83,445 of 2,374,616 requested visits, reducing fulfillment from 100\% to 96.49\%. Home exposure alone does not reduce fulfillment because requests can be sourced from other enabled homes; instead, the mean source distance increases by 6.9455~m relative to the matched reference. When both homes and activity destinations are exposed, fulfillment remains 96.49\%, showing that unavailable destinations, rather than unavailable source homes, determine unmet activity demand in this configuration.

\begin{table}[ht]
\centering
\footnotesize
\caption{Popular Times outcomes with the original L\'evy routine. Suppression results use 30 matched seeds; the adaptation row reports the corresponding 30-seed rerouting condition.}
\label{tbl:popular_times_disruption}
\begin{tabularx}{\linewidth}{@{}lrrrrX@{}}
\toprule
\textbf{Exposed nodes and policy} &
\shortstack{\textbf{Fulfillment}\\\textbf{(\%)}} &
\textbf{Unmet} &
\shortstack{\textbf{Travelers}\\\textbf{/capita}} &
\shortstack{\textbf{Source distance}\\\textbf{(m)}} &
\textbf{Interpretation} \\
\midrule
None
    & 100.00 & 0 & 17.9997 & 81.9953 & Reference. \\
Activity destinations
    & 96.49 & 83,445 & 17.3672 & 76.2669 & Requests to unavailable destinations suppressed. \\
Homes
    & 100.00 & 0 & 17.9997 & 88.9408 & Complete demand sourced from farther homes. \\
Homes and destinations
    & 96.49 & 83,445 & 17.3672 & 76.8287 & Destination loss determines fulfillment. \\
Homes and destinations; reroute
    & 100.00 & 0 & 17.9997 & 81.6733 & Unavailable destinations replaced by safe same-type POIs. \\
\bottomrule
\end{tabularx}
\end{table}

The lower source-distance values in the suppression conditions are not evidence of improved accessibility. They are calculated only for completed movements: requests whose unavailable destinations would have contributed different trips are omitted. This selection effect is why fulfillment and distance must be interpreted together.

Nearest-enabled, same-type substitution restores all 83,445 suppressed visits in both the deterministic L\'evy-off run and all 30 L\'evy-on runs. The mean displacement between the requested and receiving activity destinations is 297.89~m per rerouted traveler. Demand is redistributed among 158 receiving points of interest, and the maximum additional receiving load is 92 travelers. Compared with suppression, rerouting adds 4.8447~m per requested visit in the L\'evy-on configuration while restoring fulfillment from 96.49\% to 100\%. These outcomes quantify the cost and concentration produced by the implemented nearest-safe rule; they do not validate that residents would select the same alternatives.

Enabling the original L\'evy routine changes the home populations available when Popular Times requests are sourced, but it does not change fulfillment, travelers per capita, or occupancy per requested visit in any exposure stratum. Thus, in this matrix, the scheduled activity policy determines whether demand is fulfilled, while the L\'evy routine primarily changes the measured sourcing distance.

\subsubsection{Flood-Aware Work and School Mobility}

L\'evy Walk V2 exposes an effect hidden by the original routine: work and school attendance can be requested and audited independently of activity visits. Table~\ref{tbl:levy_v2_disruption} shows that destination exposure reduces commute fulfillment to 93.54\%, while home exposure reduces it to 94.06\%. Exposing all relevant node types combines these losses and reduces commute fulfillment to 89.15\%.

\begin{table}[ht]
\centering
\footnotesize
\caption{Popular Times V2 and L\'evy Walk V2 outcomes over 30 matched seeds. Brackets contain 95\% $t$ intervals.}
\label{tbl:levy_v2_disruption}
\begin{tabularx}{\linewidth}{@{}lrrrX@{}}
\toprule
\textbf{Condition} &
\shortstack{\textbf{Commute}\\\textbf{fulfillment}} &
\shortstack{\textbf{Commute}\\\textbf{unmet}} &
\shortstack{\textbf{Mean distance}\\\textbf{(m)}} &
\textbf{Activity fulfillment} \\
\midrule
No exposure
    & 1.0000 [1.0000--1.0000] & 2.0 & 1,393.53 & 1.0000 \\
Destinations
    & 0.9354 [0.9352--0.9355] & 352,948.0 & 1,343.80 & 0.9649 \\
Homes
    & 0.9406 [0.9404--0.9408] & 324,290.7 & 1,390.89 & 1.0000 \\
All node types
    & 0.8915 [0.8912--0.8917] & 592,460.8 & 1,335.82 & 0.9649 \\
All node types; activity reroute
    & 0.8905 [0.8903--0.8908] & 597,810.4 & 1,335.82 & 1.0000 \\
\bottomrule
\end{tabularx}
\end{table}

The shorter mean commute distances under destination and combined exposure again reflect the set of trips that remain feasible, not a travel benefit. Relative to the reference, destination exposure leaves a mean of 352,946 additional people without a completed commute and removes 49.74~m from the mean distance among completed trips. Combined exposure leaves 592,459 additional people without a completed commute and reduces the completed-trip mean by 57.71~m.

Activity rerouting in the combined condition restores Popular Times fulfillment to 100\% but does not restore work or school attendance. Commute fulfillment remains approximately 89\% and is slightly lower than in the matching activity-suppression condition. This composed result is consistent with the two routines drawing from the same represented population while retaining different destination semantics: the activity plugin can substitute another point of interest, while the commute plugin continues to report the unavailable workplace or school as unmet.

\subsubsection{Shelter Reference and Reallocation Results}

Table~\ref{tbl:shelter_reference_results} presents the principal shelter references. At initialization, 10,963 of the 12,519 configured places are occupied, leaving 1,556 immediately free. The reference evacuation fills the remaining capacity and ends with 12,519 sheltered people, or 6.84\% of the 182,936 mapped exposed population. The remaining 170,417 people retain an \texttt{in\_danger} status. Coverage is identical across seeds because the fixed capacity is fully used.

\begin{table}[ht]
\centering
\footnotesize
\caption{Mean shelter outcomes over five seeds. Waiting is accumulated over hourly simulation steps and is therefore reported in person-hours.}
\label{tbl:shelter_reference_results}
\begin{tabularx}{\linewidth}{@{}lrrrrX@{}}
\toprule
\textbf{Policy} &
\shortstack{\textbf{Coverage}\\\textbf{(\%)}} &
\shortstack{\textbf{Unresolved}\\\textbf{demand}} &
\shortstack{\textbf{Peak}\\\textbf{waitlist}} &
\shortstack{\textbf{Waiting}\\\textbf{(million person-h)}} &
\textbf{Mean distance} \\
\midrule
Fixed-initial demand; no reallocation
    & 6.843 & 170,417 & 24,195 & 3.998 & 5.511~km \\
Fixed-initial demand; reallocation
    & 6.843 & 170,417 & 24,195 & 3.950 & 5.537~km \\
Remaining-demand basis; reallocation
    & 6.843 & 170,417 & 21,227 & 3.523 & 5.607~km \\
\bottomrule
\end{tabularx}
\end{table}

Reallocation does not change final coverage or peak waitlist in the reference because the network has no unused capacity once all shelters are full. It does reduce accumulated waiting by 47,813 person-hours, a 1.20\% reduction, while increasing mean travel distance by 0.0267~km. Basing each daily request on remaining rather than initial exposure further reduces the peak waitlist by 2,968 and accumulated waiting by 427,240 person-hours relative to the reallocation reference. This change reduces repeated demand pressure but cannot increase final protection after the same capacity has been exhausted.

\subsubsection{Shelter Capacity, Exposure, and Response Rate}

Table~\ref{tbl:shelter_capacity_response} separates the quantity of shelter space from the rate at which exposed residents attempt to use it. Halving capacity reduces final coverage to 3.42\%. Doubling it raises coverage to 13.69\%, and multiplying it by five raises coverage to approximately 20.07\%. Increasing capacity from five to fifteen times the reference produces essentially no additional protection within 14 days. At this point, capacity is no longer binding: the 1\%-per-day fixed-initial response schedule limits how much of the exposed population reaches shelters during the simulated horizon.

\begin{table}[ht]
\centering
\footnotesize
\caption{Shelter capacity and daily-response sensitivity over five seeds. The reference uses 1$\times$ capacity and a 1\% daily response.}
\label{tbl:shelter_capacity_response}
\begin{tabularx}{\linewidth}{@{}lrrrrX@{}}
\toprule
\textbf{Condition} &
\shortstack{\textbf{Coverage}\\\textbf{(\%)}} &
\textbf{Unresolved} &
\shortstack{\textbf{Peak}\\\textbf{waitlist}} &
\shortstack{\textbf{Waiting}\\\textbf{(million person-h)}} &
\textbf{Binding effect} \\
\midrule
0.5$\times$ capacity & 3.423 & 176,674 & 25,499 & 4.362 & Capacity \\
1$\times$ capacity  & 6.843 & 170,417 & 24,195 & 3.950 & Capacity \\
2$\times$ capacity  & 13.687 & 157,898 & 11,676 & 0.981 & Capacity \\
5$\times$ capacity  & 20.069 & 146,223 & 1,843 & 0.033 & Response horizon \\
15$\times$ capacity & 20.069 & 146,222 & 1,841 & 0.026 & Response horizon \\
\addlinespace
0.5\% response & 6.843 & 170,417 & 11,436 & 1.755 & Slower arrival \\
2\% response   & 6.843 & 170,417 & 49,633 & 8.357 & Faster queuing \\
5\% response   & 6.843 & 170,417 & 124,364 & 21.429 & Faster queuing \\
\bottomrule
\end{tabularx}
\end{table}

Changing the response rate alone does not change final coverage because even the 0.5\% condition eventually fills the reference capacity. It substantially changes congestion: increasing the rate from 0.5\% to 5\% raises the peak waitlist from 11,436 to 124,364 and accumulated waiting from 1.75 million to 21.43 million person-hours. Faster requests are therefore not equivalent to faster protection when admission capacity is fixed.

Exposure scaling produces the corresponding demand-side effect. With 25\% of the reference exposure, coverage reaches 26.85\%, unresolved demand falls to 34,111, and accumulated waiting is 0.643 million person-hours. At 50\% exposure, coverage is 13.42\%, unresolved demand is 80,742, and waiting is 1.721 million person-hours. The 100\% reference returns to 6.84\% coverage and 3.950 million waiting person-hours.

Initial occupancy changes who can be newly admitted and how quickly queues develop, even when every scenario ends with the same 12,519 sheltered people. Empty, half-capacity, and configured-occupancy shelters produce peak waitlists of 13,339, 19,559, and 24,195, respectively. Their accumulated waiting rises from 1.250 to 2.604 and 3.950 million person-hours. The empty condition also increases mean travel distance to 6.124~km, compared with 5.537~km under configured occupancy, because a different set of exposed residents uses the newly available places.

\subsubsection{Shelter Distance and Failure Topology}

The L\'evy-scale sensitivity in Table~\ref{tbl:shelter_spatial_failures} changes where evacuees are sent without changing aggregate capacity. Coverage, unresolved demand, and peak waitlist are identical at 500, 1,000, and 3,000~m. Mean travel distance increases from 5.277 to 6.394~km across this range, and the 95th-percentile distance increases from 17.577 to 18.382~km. Thus, destination sampling changes the geographic burden while the fixed admission limit continues to determine protection.

\begin{table}[ht]
\centering
\footnotesize
\caption{Mean shelter distance and failure outcomes over five seeds. Failure rows use waitlist reallocation.}
\label{tbl:shelter_spatial_failures}
\begin{tabularx}{\linewidth}{@{}>{\raggedright\arraybackslash}Xrrrrl@{}}
\toprule
\textbf{Condition} &
\shortstack{\textbf{Coverage}\\\textbf{(\%)}} &
\textbf{Unresolved} &
\shortstack{\textbf{Mean distance}\\\textbf{(km)}} &
\shortstack{\textbf{P95 distance}\\\textbf{(km)}} &
\textbf{Capacity removed} \\
\midrule
L\'evy scale 500~m   & 6.843 & 170,417 & 5.277 & 17.577 & None \\
L\'evy scale 1,000~m & 6.843 & 170,417 & 5.537 & 17.699 & None \\
L\'evy scale 3,000~m & 6.843 & 170,417 & 6.394 & 18.382 & None \\
\addlinespace
No shelter failure        & 6.843 & 170,417 & 5.537 & 17.699 & None \\
Largest shelter disabled  & 6.573 & 170,912 & 5.522 & 17.714 & 495 places \\
Highest-capacity 10\% of shelters disabled & 4.675 & 174,384 & 5.731 & 17.917 & 3,967 places \\
Humait\'a regional failure & 6.843 & 170,417 & 5.537 & 17.699 & 0 places \\
\bottomrule
\end{tabularx}
\end{table}

Disabling the single largest shelter removes 495 places and reduces coverage by 0.271 percentage points. Disabling the 11 highest-capacity shelters---10\% of the 104-node network after rounding---removes 3,967 places and reduces coverage by 2.169 percentage points. The configured Humait\'a failure selects no shelter in the stored input and consequently removes no capacity or changes no outcome. This null result is informative about scenario topology: naming an affected region has no shelter-service effect when the modeled network contains no shelter in that region.

Reallocation again changes waiting rather than final coverage in the failure scenarios. It reduces accumulated waiting by 46,546 person-hours after the largest-shelter failure and by 35,265 person-hours after failure of the 11 highest-capacity shelters, while adding approximately 0.027 and 0.030~km to mean travel distance, respectively. Once every remaining enabled shelter is full, no reallocation policy can recover the protection lost with disabled capacity.

\subsection{Remarks}

The results demonstrate why environmental state and response policy must remain separate. The water-level plugin supplies the same type of availability change to multiple mobility models, yet the outcomes differ because Popular Times may suppress or substitute an activity destination, L\'evy Walk V2 records requested work and school demand as unmet, and the shelter workflow creates new protection-seeking movement subject to capacity. The framework supports these responses without embedding any one of them in the core node-state mechanism.

Destination substitution is effective only for the demand to which it is applied. The nearest-safe policy restores all configured marketplace, restaurant, and pharmacy visits, but it adds travel and receiving load and does not compensate for inaccessible workplaces or schools. Similarly, shelter reallocation uses distributed spare capacity and reduces waiting exposure, but it cannot replace failed or insufficient capacity. Coverage, unmet demand, waiting, distance, and receiving load therefore capture distinct effects and should be reported together.

The shelter sensitivity results further show that capacity and response speed are joint constraints. Under reference capacity, faster evacuation requests mainly create longer queues. Under sufficiently expanded capacity, the 14-day response schedule limits final protection. Failure consequences depend on how much capacity is removed and where shelters exist, rather than solely on the number of affected regions.

These results remain conditional on the implemented assumptions. Water thresholds, hourly activity profiles, same-type nearest substitution, commute attendance rates, exposure multipliers, shelter capacities, response rates, and failure selections are experiment inputs. Mobility uses geographic rather than transport-network distance. The shelter workflow models static initial exposure, simple capacity admission, and no demographic priority, recovery, or return-home decision. Its five-seed catalog is a screening analysis of stochastic destination sampling, not an exhaustive uncertainty analysis or a forecast of observed evacuation behavior.
